\section{From the boundary of the weight chamber to a slope}
\label{sec:flag-wall-slope}

The two tests of Section~\ref{sec:covolume-flag-test} give values of the arithmetic functional on specific polynomials. A slope estimate at the chosen weight requires information on the whole weight polynomial. To turn this information into
a slope estimate, we use the limiting behavior at the walls to
constrain a homogeneous weight polynomial. The integral
comparison on the deep cover then supplies short sections and
bounds on the moving degree. The next chapter returns the
component estimates to the original level.

Fix a marked flag space of rank \(N\) as in Section~\ref{sec:covolume-flag-test}. At each inserted prime, compute the mixed arithmetic test on the common integral model of Lemma~\ref{lem:flag-finite-frobenius-model}. For its embedding set \(\mathcal T\), let \(d_F=|\mathcal T|\binom N2\) and \(n=d_F+1\). To a weight array \(u=(u_{\tau i})\) associate the class
\[
 L_u=-\sum_{\tau,i}u_{\tau i}x_{\tau i}.
\]
Let \(\Phi_B\) be the
linear functional obtained by taking the smoothing limit of the
normalized arithmetic degree at cutoff \(B\ge2\). We determine
\(\Phi_{B_0}\) on forms of degree \(n\), for a fixed real
\(B_0\ge2\), by varying an auxiliary integer cutoff \(B\).

In particular the weight can later determine \(B_0\).
The arithmetic and reference functionals on the same array of weights are
\begin{equation}\label{eq:wall-weight-polynomials}
 H_B(u)=\Phi_B(L_u^n),\qquad H(u)=h_{\rm fl}(L_u^n).
\end{equation}
The geometric top-degree polynomial is
\[
 G(u)=\deg_{\rm fl}(L_u^{d_F})
   =d_F!\prod_{\tau}\prod_{a<b}\frac{u_{\tau a}-u_{\tau b}}{b-a}.
\]
The last equality follows by expanding the power and taking
the signed top coefficient. In the resulting determinant the
factorials in the multinomial coefficient give
\(\prod_{i=0}^{N-1} i!=\prod_{a<b}(b-a)\) per embedding.
Thus \(G\) is positive in the open dominant chamber, where
the entries at each embedding decrease strictly.

\subsection{Heights on the marked flag model}
\label{subsec:flag-point-heights}

The behavior near a wall will be inferred from heights on the compactified flag model. The height bounds must include the boundary as well as the open abelian locus.

Fix the deep cover at a nonterminal letter of either run, with CM
field \(E\), real degree \(f\), rank \(N\), and signature \((j,N-j)\).
Write \(A\) for its auxiliary ample metrized line and \(h_A\) for
the associated absolute point height. The perturbed line \(M_B(v,s)\)
has a summand \(sA\). Its root part has weight
\[
 w=v-sd_0,
\]
where
\(d_0\) is fixed on the cover. The arithmetic degree of a lattice over
a number field will always be divided by the degree of the field.
In this subsection constants may depend on the fixed cover. We shall
use a threshold in \(s\) independent of \(B\), of the point and its
field, and of the weight in a bounded set. This order of quantifiers
will allow us to take a limit for each cover before applying a
separate uniform estimate as the covers vary.

\paragraph{Lie lattices.}

The first height contribution is the degree of a line in the integral Lie algebra.

Let \(\mathcal L\) be the integral Lie algebra of a semiabelian model
over a number ring with principally polarized abelian generic fibre of
dimension \(g\). Use the polarization metric \(\pi H_\sigma\) on the
Lie algebra at each embedding, where \(H_\sigma\) is the Hermitian
form of complex uniformization. Put
\[
 h_*=-\widehat{\deg}\mathcal L
\]

and \(H=\max(h_*,0)\). For \(g\le g_0\), every integral sublattice
\(U\subset\mathcal L\) satisfies
\begin{equation}\label{eq:flag-lie-slope}
  \widehat{\deg} U\le \rk U\,C(g_0)\log(2+H).
\end{equation}
To see the logarithmic dependence, take a nonzero algebraic tangent
direction \(x\) and let \(a\) be the degree of the saturated line it
spans. Saturation of this line commutes with finite extension of the
number field, by flatness and the torsion-free quotient. Both its
normalized degree and its metric therefore persist over the fields
used below.

For an integer \(n\ge2\), put \(m=4n\) and let \(L\) be the
\(m\)-th power of the symmetric theta line, with its rigidification.
Over a finite extension take the semistable Moret--Bailly model
of \((A,L)\). The line has a cubical extension and the polarization
kernel of \(L^{\otimes 2}\) extends as a finite flat subgroup.
On this model Bost's semistable direct-image theorem gives a semistable lattice of sections with the cubical probability \(L^2\) metric. Its rank is 
\[
 \operatorname{rk}H^0(A,L)=m^g=(4n)^g.
\]
 In the Faltings normalization
of Gaudron's 2019 statement its slope is
\[
 -h_F/2+(g/4)\log(m/(2\pi^2)).
\]
For a principal polarization,
\[
 h_* = h_F+(g/2)\log\pi,
\]
because the tangent metric here is
\(\pi H_\sigma\). Thus the slope in these norms is
\[
 \nu=-h_*/2+(g/4)\log(m/(2\pi)).
\]

\cite{bost-mb-flag}\label{cor:flag-point-heights-1}

\begin{proposition}[Point-height bounds on the selected flag]
\label{prop:selected-flag-point-heights}
For the marked model the relative Lie lattices satisfy the
slope bound \eqref{eq:flag-lie-slope}. The resulting interior
and boundary point-height bounds use the same root metric.
On every component there is a dense set on which the two inequalities hold simultaneously
\[
 |h_A|\le C,\qquad |h_0(w)|\le C\|w\|_\infty\quad\text{for all }w.
\]
 The set is fixed before
\(B,v,s\).
\end{proposition}

\begin{proof}
\proofstep{Filtering by the evaluations at successively chosen points.}
Filter the space of sections by vanishing of the successive
directional derivatives at zero. The first \(n+1\) quotients have
rank one. Indeed \(\Theta^4\) separates the tangent at zero. Take two
sections of \(\Theta^4\), one nonzero at zero and the other with order
one along \(x\). Call the two sections \(s_0,s_1\). The powers satisfy
\[
 s_0^{\,n-i}s_1^i\in H^0(A,\Theta^{4n}),\qquad
          \operatorname{ord}_{x,0}(s_0^{\,n-i}s_1^i)=i,
                  \quad 0\le i\le n .
\]
Here the order is along the chosen tangent \(x\) at zero.
They exhibit a nonzero class in each of the successive
one-dimensional jet quotients.
Evaluate the \(i\)-th derivative on the kernel of the preceding
jets using the rigidification.

The derivative is integral on an integral tangent and an integral section. In the group law this derivative is the
coefficient of \(e_1\cdots e_i\) at the sum of \(i\) copies of that
tangent, in the algebra with \(e_k^2=0\). A change of trivialization
uses only its value at zero on this kernel. Writing \(j_i(s;x)\) for this derivative of a section \(s\) in the tangent \(x\), homogeneity of the multilinear coefficient gives
\[
 j_i(s;bx)=b^i j_i(s;x).
\]
Thus its finite-place image
pairs integrally with the \(i\)-th power of the saturated tangent
line.

At an embedding let \(\rho\) be the shortest nonzero period in
the \(H_\sigma\)-norm. The archimedean operator bound for this jet,
after removing \(\|x\|_{\pi H_\sigma}^i\), is
\[
  \beta_i(\rho)=(e m i)^{i/2}\sqrt{e\,g!}\,
       \max\{(2/(\sqrt\pi\rho))^g,m^{g/2}\}.
\]
Use \(1\) for the first factor when \(i=0\). In a translated
Appell--Humbert frame the centred Gaussian is
\(\exp(-\pi m H_\sigma(z,z))\) for the squared norm. A ball of
radius
\[
 \min(\rho/2,1/\sqrt{\pi m})
\]
injects up to its boundary
into the torus, and the Gaussian on it is at least \(e^{-1}\).

The holomorphic mean-value inequality and the volume of this ball
give the last two factors in the comparison of the supremum and
probability \(L^2\) norms. The fundamental period parallelepiped
has volume one in these coordinates by principality. Cauchy's
inequality along the unit tangent with radius
\(\sqrt{i/(\pi m)}\) in the \(H_\sigma\)-norm, and \(i!\le i^i\),
give the first factor with the stated tangent metric. The calculation
applies to every complex linear combination of sections in the
kernel. \label{cor:flag-point-heights-2}

\proofstep{The quotient heights and the product formula.}
The determinant of a rank-one quotient and the product formula
now give its normalized degree at most
\[
 -ia+\operatorname{mean}_\sigma\log\beta_i(\rho_\sigma).
\]

The remaining kernel has degree at most its rank times \(\nu\) by
semistability. Subtracting this from the full degree gives
\[
 (n+1)\nu\le -\tfrac12n(n+1)a+
           \sum_{i=0}^n\operatorname{mean}_\sigma\log\beta_i(\rho_\sigma).
\]
For principal polarization the matrix lemma bounds the mean of
\(\rho_\sigma^{-2}\) by \(14\max(1,h_*-(g/2)\log\pi,\log(g!))\).
Apply concavity of \(\log\) to \(1+\rho_\sigma^{-2}\). The mean in
the last display is at most
\[
  \tfrac i2\log(e m n)+\tfrac g2\log m+
        \tfrac g2\log(2+H)+O_g(1).
\]
Field extension repeats each base embedding equally often in this
mean.

Take \(n=2\lfloor H\rfloor+2\). Then \(H/n\le1\) and
\(\log n\le\log(2+H)+\log 2\). Dividing by \(n(n+1)/2\) gives
\(a\le\log(2+H)+O_g(1)\). Let \(U\) be a Hermitian sublattice of positive rank. The absolute
Siegel lemma gives, after finite extension, a line \(L\subset U\)
with
\[
    \widehat{\deg} L\ge \frac{\widehat{\deg} U}{\operatorname{rk}U}
                         -\tfrac12\log\operatorname{rk}U-1.
\]
The inequality extends to \(L\) from its saturation \(L^{\rm sat}\) by
\[
 \widehat{\deg} L\le\widehat{\deg} L^{\rm sat},\qquad \operatorname{rk}L=\operatorname{rk}L^{\rm sat}.
\]
This proves \eqref{eq:flag-lie-slope}, taking
a maximum over the finitely many dimensions up to \(g_0\). Rank zero
contributes zero. \cite{matrix-flag,absolute-siegel-flag}

\paragraph{Interior points.}
Take the semiabelian integral model of the abelian variety at the
given algebraic point after semistable extension. Its dimension is
\(fN\). Pull the marked flag back with the same point, intersect the
rational steps with the Lie lattice and give them the induced
metrics. If \(W_k\) denotes this intersection, the determinant comparison reads
\[
 h_{\rm step}=-\widehat{\deg}(\det W_k)+O(1).
\]
Indeed the product of the successive local quotient generators is
a determinant generator, and the product of their quotient norms is
the Gram determinant norm. At the finitely many primes where
integral eigenprojections require a denominator the field
discriminant bounds the change in the normalized valuation.

At infinity the polarized metric pulled back at the point
is the Hodge Gram form in its marked frame. The successive Gram
determinant ratios therefore agree at every embedding. If a rational determinant frame over \(K\) is rescaled by \(a\in K^\times\), the changes cancel by
\[
 \sum_v[K_v:\mathbb Q_v]\log |a|_v=0,
\]
with absolute values extending the standard ones. Applying the determinant comparison to the full Lie lattice gives
\(h_F=h_*+O(1)\). If \(h_{\rm step}\) denotes any one of the step
heights, \eqref{eq:flag-lie-slope} consequently gives
\[
                    h_{\rm step}\ge-C\log(2+\max(h_F,0))-C .
\]
Applying it to the entire lattice bounds the negative part of
\(h_F\) by the same quantity in absolute value.
\label{cor:flag-point-heights-3}

\proofstep{Root heights in the interior metric.}
At \(\tau\) write \(R_{\tau i}\), \(0\le i<N\), for the
graded root heights in the Hodge metric and
\[
 P_{\tau k}= \sum_{i<k} R_{\tau i}.
\]
The first \(j\) entries are the
graded Lie heights at \(\tau\). The remaining entries are
the negatives of those at \(\bar\tau\) in reverse order.
To bound the full sum \(P_{\tau N}\) in absolute value, use that
the full determinant and its inverse have integral sections up
to fixed denominators, nonzero at all generic-fibre points,
on the determinant cover. The smooth reference metrics bound their
norms above. Their changes to the Hodge metric for this central
weight are bounded, since its potential extends smoothly. Thus the step height in the first block reads
\[
 h_{\rm step}(\tau,k)=-P_{\tau k}\qquad(k\le j).
\]

For \(k>j\) it is the step height of dimension \(N-k\) at
\(\bar\tau\) minus \(P_{\tau N}\). Both have the preceding lower
bound. Moreover the sum of the two unsigned-block Lie determinants
in the Hodge metric, with negative sign, reads \(h_F+O(1)\) after
summing over \(\tau\). The two block sums differ by \(P_{\tau N}\).
Consequently
\[
  \sum_\tau -P_{\tau j}=\tfrac12 h_F+O(1).
\]
The field and level are fixed in these \(O(1)\) terms.
\label{cor:flag-point-heights-4}

Let \(\delta_{\tau k}=w_{\tau,k-1}-w_{\tau k}\) for \(1\le k<N\).
For a dominant \(w\) of norm at most \(2\), and
\(0\le b\le\min_\tau\delta_{\tau j}\), summation by parts gives
\[
 h_{\rm H}(w)=
 \sum_\tau\Big(\sum_{k=1}^{N-1}\delta_{\tau k}(-P_{\tau k})
                      -w_{\tau,N-1}P_{\tau N}\Big)
 \ \ge\ \tfrac b2\max(h_F,0)-C\log(2+\max(h_F,0))-C .
\]
To get the inequality, reserve \(b\) at each middle gap. The
remaining coefficients are nonnegative and their sum at each
embedding is at most \(4\). This also shows why the possibly
unbounded individual middle heights must first be added with a
common coefficient.

\paragraph{Boundary points.}
Use the smooth reference metric at a boundary point and the actual
extension by a torus of its polarized abelian part. After extending
the residue number field one obtains its Lie lattice \(L\), a map
\(q:L\to L'\) to the lattice of the abelian part, and a frame
\(\epsilon\) of the torus. The rational map is onto and its integral
kernel is the span of \(\epsilon\). On these data there is a fixed
\(D\ge1\) such that
\[
   {\rm Gram}(\epsilon)\ \ge D^{-1}I,\qquad
   {\rm Gram}(qv)\ \le D\,{\rm Gram}(v)
\]
for each tuple \(v\), at every embedding.

For the first comparison
split the torus over a finite cover of a dense open of each
boundary stratum. Its cocharacter basis is saturated in the
integral Lie algebra by the integral torus inclusion. Its norm
comparison on the split stratum is the torus-frame bound. Apply
Noetherian induction to the closed complement, keeping the maximum
of the finitely many bounds. The second comparison comes from the
local Hermitian bound on the map to the polarized quotient and a
finite cover of the compact complex charts. Both comparisons are
stable after extension to the common field of definition of the
point and frame. \label{cor:flag-point-heights-5}

\proofstep{The Lie quotient at the boundary.}
For \(W\subset L\) put \(W_0=W\cap\ker q\). Its coordinate lattice
in the torus basis has degree at most zero for the standard
metric, by a nonzero integral minor and the product formula.
Consequently
\[
 \widehat{\deg}W_0\le(\rk W_0)\log D/2.
\]

Using the quotient metric on \(W/W_0\), the second Gram inequality
gives
\[
 \widehat{\deg}W\le \widehat{\deg}W_0+\widehat{\deg}(qW)
                            +\tfrac12\rk(qW)\log D .
\]
In particular
\[
 h_*'=-\widehat{\deg}L'\le -\widehat{\deg}L+O(1),
\]
by taking \(W=L\) and using
\(qL\subset L'\) of full rank. The smooth determinant height
\(h_\omega=-\widehat{\deg}L+O(1)\) satisfies
\[
 |h_\omega|\le C(1+|h_A|)
\]
by ample height comparison.
Also \(h_A\ge-C\), using a finite generating family in an ample
power. Applying \eqref{eq:flag-lie-slope} to \(qW\) gives the
lower bound \(-C\log(2+C+h_A)-C\) for each smooth step height,
using the determinant comparison for the saturated marked steps
of \(L\). The argument includes a zero-dimensional abelian part
by taking \(qW=0\).

The smooth root heights have the same signed reversal as the
interior heights and bounded full sum at each \(\tau\). For
\(k>j\) use the conjugate step of length \(N-k\), which contributes
the same logarithmic loss. Summation by parts for dominant \(w\)
of norm at most \(2\) therefore gives
\[
 h_0(w)\ge -C\log(2+C+h_A)-C .
\]
Here \(h_0(w)\) uses the empty finite-prime set and smooth metrics.
\label{cor:flag-point-heights-6}

\paragraph{Upper heights on a dense set.}

The lower height bounds apply to arbitrary points. To bound the intersection in the other direction we need upper bounds on points that are dense in each component.

Fix a component, embed it by a power of the ample line, and
choose an affine chart on which the projective coordinate ratios
\(z_1,\ldots,z_m\) are regular. Noether normalization gives a
finite map
\[
 \pi\colon U\longrightarrow\mathbb A^d_K,\qquad
 K[t_1,\ldots,t_d]\hookrightarrow K[U],\qquad
 K(U)/K(t_1,\ldots,t_d)\text{ finite}.
\]
For \(0\ne a\in K[U]\), its field norm is nonzero.
Clear denominators and restrict to the open set over which the
norm specializes as the determinant of multiplication by \(a\).
The tuples of roots of unity are Zariski dense in \(\mathbb A^d\),
so take one in this open set on which the norm is nonzero. At that
tuple the determinant is invertible; every point of its nonempty
fibre lies in \(D(a)\). The union of the fibres over torsion tuples
therefore meets every nonempty principal open of \(U\).

Each \(z=z_j\) satisfies a monic equation
\(z^s+\sum_{i<s}f_i(t)z^i=0\) with \(f_i\in K[t]\).
Write \(f_i=\sum_\alpha c_{i\alpha}t^\alpha\). For a torsion tuple
and an absolute value above \(v\) the monic root bound gives
\[
 \max(1,|z|_v)\le
 \begin{cases}
  1+\sum_{i,\alpha}|c_{i\alpha}|_v,& v\mid\infty,\\
  \max(1,\max_{i,\alpha}|c_{i\alpha}|_v),& v\nmid\infty .
 \end{cases}
\]
Let \(C_v\) be the maximum of these bounds over the finitely
many coordinate ratios. It equals one outside a fixed finite
set. With the normalized local weights \(n_v\) of \(K\), summing
over the places of any point field above \(v\) gives
\begin{equation}\label{eq:root-fibre-height}
       h([1:z_1:\cdots:z_m])\le\sum_v n_v\log C_v.
\end{equation}
Passing from the projective height in \eqref{eq:root-fibre-height}
to the fixed ample reference metric changes the bound by \(O(1)\).
Lift these closed points to marked algebraic points using the
component's point map. Its image contains every such closed point,
which is the density test on the marked component.

On the same set every smooth root height \(h_x\) is bounded in
absolute value. On the common root model, twist the root and its
dual by sufficiently high powers of the pulled-back ample line so
that the resulting sections span on the generic fibre. Clear fixed denominators.
For each point use a section nonzero there. Bounded norms of the
finite list give both \(\pm h_x\le C_x(|h_A|+1)\).
Thus one dense set per component satisfies
\(|h_A|\le C\) and
\[
 |h_0(w)|\le C\|w\|_\infty
\]
simultaneously for all \(w\). The same dense set can thus be used as \(B,v,s\) vary in the subsequent intersection bounds.
\label{cor:flag-point-heights-7}
\end{proof}

\subsection{The perturbed metric and the vertical curves}
\label{subsec:flag-wall-positivity}

Height estimates over number fields control horizontal points. For the arithmetic intersection near a wall we must also control curvature and degrees on the vertical curves.

For the same marked point we compare the smooth reference
height \(h_0(w)\), the interior Hodge height \(h_{\rm H}(w)\), and
the height \(h_M\) of \(M_B(v,s)\) before taking any
intersection. With \(M'\) the marked modulus, the finite set of primes at which the metric changes is
\[
 Q_B=\{p:p\text{ prime},\ p\le B\}
      \mathbin{\cup}\{p:p\text{ prime},\ p\mid 2|\operatorname{disc}E|M'\}.
\] Use the Hodge--Tate lattice at \(p\in Q_B\)
and the de Rham lattice elsewhere.

\paragraph{Smoothing and the finite change of lattice.}
On the complex toroidal compactification use the Kähler reference
form \(\Omega\) and the Hodge potential relative to its smooth
root metric. For a proper initial flag step this is the difference
of the logarithmic determinant metrics. For the relative potential \(u\) and the smooth reference curvature \(\theta\), positivity on the open part reads
\[
 \theta+dd^c u\ge0.
\]
Indeed these fundamental lines are the dual determinants of the first
steps in the Hodge flag, with the conjugate flag in reverse order in
the second block. In period coordinates the curvature on the base
is the Hodge form on a subspace, and variation of the subspace adds
the squared norm of its second fundamental form to the dual
determinant curvature. Past the middle use the conjugate step and
the flat full determinant.

In the toroidal logarithmic charts the
Hodge metric and its inverse have polynomial growth in the boundary
radius
\[
 r=e-\log\|s_D\|^2,
\]
where \(s_D\) cuts out the boundary.
Taking the logarithm of a determinant bounds the potential by
\(C\log r+C\) in absolute value. A locally plurisubharmonic
potential with this sublogarithmic growth extends across a normal
crossing divisor. To see this, add \(\delta\sum\log|z_i|\) at its local branches,
extend across them by the upper limit and take the upper
semicontinuous regularization as \(\delta\downarrow0\).
The resulting potential is bounded above on each compact chart.
For the full central weight use the smooth extension of the
flat determinant potential, whose curvature cancels its reference
curvature identically. Its coefficient may have either sign.
\label{cor:flag-wall-positivity-1}

\begin{proposition}[A fibrewise nef perturbed root metric]
\label{prop:perturbed-flag-metric}
With the permitted smoothing and finite cutoff, \(M_B(v,s)\)
on each of the two selected marked covers has an archimedean semipositive metric. On every vertical curve \(C\) its degree satisfies
\[
 \deg_C(M_B(v,s)|_C)\ge0.
\]
 It satisfies the interior, boundary and dense-set
height comparisons \eqref{eq:flag-actual-ht}.
\end{proposition}

\begin{proof}
\proofstep{Regularizing the fundamental potentials.}
Take \(L_s=s^{-2}\) and smooth the proper fundamental potentials by
the holomorphic part of the Chern exponential map at radius
\(\eta=\exp(-L_s)\). In Demailly's regularization theorem the
background form is closed, its sum with \(dd^c\) of the potential
is a positive current, and the compact Kähler metric admits a
curvature lower bound on orthogonal tangent pairs. A multiple of
\(\Omega\) provides that bound. The regularized potential
\(\psi_\eta=\rho_\eta\psi+K\eta^2\) majorizes \(\psi\) on the
open part. It satisfies, uniformly on the compactification,
\[
 -C(1+\log L_s)\le\psi_\eta\le C,\qquad
 \theta+dd^c\psi_\eta\ge -C\frac{\log L_s}{L_s}\Omega .
\]
Here are the rates in this use of regularization. In each crossing
chart the ball mean of \(-\log |z_i|\) truncated below at zero is
\(O(1+\log(1/\eta))\).

To obtain this bound, rescale the ball when its centre is within
twice the radius of the hyperplane, and use the distance bound
on the other balls. The normalized kernel is bounded by
\(C\eta^{-2d}\) on a ball of radius \(C\eta\) in complex
dimension \(d\). The smooth chart factors change the estimate by
a constant. Jensen's inequality for the logarithm of the boundary
radius gives the lower bound on \(\psi_\eta\). For the curvature
bound Demailly's horizontal Hessian estimate costs
\[
 O(\eta^2+\partial_{\log\eta}\psi_\eta)\Omega.
\]
The radial
Hessian bound gives convexity in \(\log\eta\) after a further
fixed multiple of \(\eta^2\) is added. Comparing with a fixed
small radius gives a secant bound \(O((1+\log L_s)/L_s)\) for
the derivative, and gives the displayed estimate for small \(s\).
There are finitely many fundamental steps and charts for the fixed
model. \cite{demailly-flag}\label{cor:flag-wall-positivity-2}

Form the weighted sum of these potentials by the successive gaps
of \(w\), and its last coefficient for the full determinant. Call
the result \(\phi_s(w)\). In each block of consecutive proper steps the coefficients sum by telescoping, 
\[
 \sum_{i=a}^{b-1}(w_i-w_{i+1})=w_a-w_b\le2\|w\|_\infty.
\]
 They are nonnegative by dominance, so their total on the fixed model is \(O(\|w\|_\infty)\).
The smooth auxiliary line \(A\) has curvature at least \(c\Omega\),
with \(c>0\) on every component. For \(\|w\|_\infty\le M\), choose
\(s_0>0\) so that \(CMs\log(1/s)\le c\) for \(0<s\le s_0\).
The curvature of the perturbed metric then satisfies
\[
 \operatorname{curv} M_B(v,s)
 \ge \bigl(sc-C\|w\|_\infty s^2\log(1/s)\bigr)\Omega
 \ge s\bigl(c-CMs\log(1/s)\bigr)\Omega\ge0 .
\] Its potential is at least the open Hodge
potential on the interior, globally bounded above by a constant,
and globally bounded below by \(-C\log(1/s)-C\).

\proofstep{Finite changes of the root metrics.}
For each root, at a changed prime \(p\) the normalized point-height
difference is at most
\[
   J_p=\frac{\log p}{p-1}+
     \frac{v_p(|\operatorname{disc}E|)}{[E:\Q]}\log p .
\]
Indeed the Hodge--Tate and de Rham integral steps differ by the
Hodge--Tate lattice loss \(1/(p-1)\) and the normalized different
in the labelled projection. The same estimate holds on the reversed
dual roots. In summing local readings on the field of a point,
the degrees above \(p\) add to the degree of that field. The second
sum in \(J_p\) is bounded by \(\log\operatorname{rd}_E\), and
\[
 \sum_{p\le B}\log p/(p-1)=O(1+\log B)
\]
by the elementary
Chebyshev bound and partial summation. The extra marked primes have
fixed total cost. Hence switching the finite lattices from the
empty set to \(Q_B\) costs \(O(\|w\|_\infty\log B)\), uniformly
over the points. \cite{scholze-weinstein-flag}\label{cor:flag-wall-positivity-3}

\proofstep{Assembling the perturbed point height.}
At a point the arithmetic height is additive. Its terms are the
smooth height at \(Q_B\), the average of \(\phi_s(w)\) over the
complex readings of the point, and \(s h_A\). The preceding
comparisons give
\begin{align*}
  h_M&\ge h_{\rm H}(w)+s h_A-C\log B &&\text{in the interior},\\
  h_M&\ge h_0(w)+s h_A-C\log(B/s) &&\text{everywhere},\\
  h_M&\le h_0(w)+s h_A+C\log B &&\text{everywhere}.
\end{align*}
All three use the same finite lattices. Given \(\|v\|_\infty\le1\),
let the middle gaps of \(v\) be at least \(t\), and the other gaps
at least a fixed \(c_0>0\).

Take \(s=t^2\),
\(t\le\min(c_0,1)\) and \(4t\|d_0\|_\infty\le1\). Every
shifted adjacent gap is at least \(t/2\), since it loses at most
\(2t^2\|d_0\|_\infty\). Also \(\|w\|_\infty\le2\).
Use \(b=t/2\) in Section~\ref{subsec:flag-point-heights}. With
\(h=\max(h_F,0)\) the resulting bounds, after enlarging \(C\), are
\begin{align}
 h_M &\ge \tfrac t4 h-C\log(2+h)-C\log B-C
           &&\text{in the interior},\notag\\
 h_M &\ge t^2 h_A-C\log(2+C+h_A)-C\log(B/t^2)-C
           &&\text{on the boundary},\label{eq:flag-actual-ht}\\
 h_M &\le C\log B+C &&\text{on a dense set in each component}.
           \notag
\end{align}
The lower bound for \(h_A\) applies also in the interior.

\paragraph{Degrees at a Hodge--Tate prime.}
For \(p\in Q_B\) pull the root lines to a common normal projective
cover splitting level \(p^2\) and the roots of unity needed for its
torsion characters. Consider a first step of length \(k\) in an
eigenflag. In each affine tangent frame the characters on the
\(p^2\)-torsion have integral representatives for their logarithmic
differentials. Evaluate their \(k\)-fold minors on the step
determinant. At an integral test point the largest minor gives
the Hodge--Tate determinant norm. If \(m_\alpha\) are the normalized minors in the induced frame, integrality and the unit minor give
\[
 \min_\alpha v(m_\alpha)=0.
\]

On overlaps
the differences of the two representatives for each minor have
normalized absolute value strictly less than one. Thus their
residue sections agree. More explicitly, the integral point
criterion on the normal model gives regularity at the special
fibre from the generic regularity and the integral readings.
At a point of that fibre take an integral test point whose
closed specialization detects units in the local ring. A unit
minor there gives a generator on a neighbourhood, and the strict
overlap inequality gives equality of the restrictions to every
reduced vertical curve. The residue sections generate the step line on each such
curve. For a curve \(C'\) above \(C\), let \(L_k\) be the
step line and let \(A_{\rm base}\) account for the discriminant
correction after clearing by \(d_E=[E:\Q]\). The line measured
by the generated minors is
\[
 L_k^{\otimes d_E}|_{C'}
\]

times the pullback of \(A_{\rm base}\). Global generation on \(C'\) gives nonnegative degree, and the
base line pulls back with degree zero. Consequently
\[
 0\le\deg_{C'}(L_k^{\otimes d_E}\otimes A_{\rm base}|_{C'})
    =d_E\deg_{C'}L_k .
\] To descend choose a closed
point of the fibre over the generic point of \(C\) and take
its closure \(C'\); the map has a positive finite degree
\(n\). The projection formula
\[
 \deg_{C'}L_k=n\deg_C L_k
\]

gives the same inequality on \(C\). This construction uses the determinant on the first \(k\)
steps, including at a prime of the discriminant.
\label{cor:flag-wall-positivity-4}

The full signed determinant at each embedding has degree zero on
these curves. The flat determinant section has a constant normalized
Hodge--Tate valuation on the component at the chosen place. The
section and its inverse therefore induce nef lines there. The
comparison uses the determinant frame at that place and its
valuation at a test point on the same component. For a step past
the split \(j\), the negative prefix is a negative prefix of the
conjugate Lie flag plus a full signed determinant. Its degree is
again nonnegative. Clear the rational weight denominators.
For an antitone weight \(w_1\ge\cdots\ge w_N\),
\[
 w=\sum_{k=1}^{N-1}(w_k-w_{k+1})\mathbf1_{[1,k]}
                         +w_N\mathbf1_{[1,N]}.
\]
The gaps are nonnegative, and the full signed determinant
has degree zero on the tested curves, so the root combination
has nonnegative degree. Adding \(sA\) preserves it. \label{cor:flag-wall-positivity-5}

\paragraph{Degrees at a de Rham prime.}
For \(p\notin Q_B\) we have \(p>2\), \(p>B\), and \(p\) divides
neither the field discriminant nor the marked level. Test on an
integral vertical curve. Under a dominating projective model a
contracted curve has zero pullback degree. In the other case
normalize a dominating curve over \(\overline{\mathbb F}_p\).
This multiplies all degrees by the same positive factor. The
unramified action splits into eigenbundles.

In using the semiabelian
family on this curve, one can pass to a finite cover, separate the
constant split torus and take the polarized abelian generic part.
For a subbundle \(W\), let \(W'\) denote its image in the quotient and \(K\) its kernel in the constant torus part. Then
\[
 \deg W=\deg W'+\deg K\le\deg W',
\]
since \(K\) is a subsheaf of a trivial bundle.
The full eigenbundle degrees are preserved under removal of the
constant torus. Further finite pullback multiplies degrees by
the map degree. Thus the flag prefix upper bounds below descend
with the same factor as the total degrees.
\label{cor:flag-wall-positivity-6}

\proofstep{The Hodge and conjugate filtrations.}
Write \(U=\bigoplus_\sigma U_\sigma\) for the Lie bundle of the
resulting family and \(D_\sigma\) for its de Rham eigenbundles.
The polarized Hodge and conjugate filtrations on the abelian
open give
\[
 \begin{aligned}
 0&\to U_{\bar\sigma}^\vee\to D_\sigma\to U_\sigma\to0,\\
 0&\to F^*U_\sigma\to D_{\varphi\sigma}
          \to F^*(U_{\bar\sigma}^\vee)\to0 ,
 \end{aligned}
\]
where \(F\) is Frobenius and \(\varphi\) is its permutation
of the embeddings. For precision, the sequence on the good open must extend as an
exact sequence of vector bundles at each missing point. In the
curve comparison one first makes the finite Frobenius pullback,
using the Frobenius twist of the base field structure on the smooth
curve. Over the abelian open take the two filtrations on primitive
de Rham cohomology with its polarization pairing.

At a missing
point, in the completed discrete valuation ring let \(F_0\) denote
the rational Hodge subspace and \(\Omega_0\) the lattice of integral
invariant differentials. The integral comparison for the
semiabelian torsion tower gives a finitely generated lattice
\(\Lambda\) stable under the endomorphisms with Rosati adjoints.
Its two conditions of adaptation for a rational subspace \(V\)
and lattice \(N\subset V\) are
\[
 \Lambda\cap V=N,\qquad \Lambda+V=N^{\vee,\beta},
\]
where the right side means the inverse image of the integral
linear dual under the polarization pairing \(\beta\).
The integral polarized comparison for the \(p\)-divisible tower
gives both identities at \((F_0,\Omega_0)\). It also gives
\[
 \Lambda\cap W=R\langle\operatorname{Fr}(\Lambda)\rangle,
 \qquad W=K\langle\operatorname{Fr}(H)\rangle ,
\]
where \(R\) is the completed local ring and \(H\) the rational
primitive de Rham space over its fraction field \(K\). This is the
integral Frobenius saturation of the transported lattice.

Since
Frobenius kills \(F_0\) and
\(\Lambda+F_0=\Omega_0^{\vee,\beta}\), this intersection is exactly
\(N=R\langle\operatorname{Fr}(\Omega_0^{\vee,\beta})\rangle\).
The same comparison sends \(\Lambda\) under Verschiebung into the
integral span of the images of \(\Omega_0\) in the Frobenius-twisted
space. In the Frobenius--Verschiebung adjunction, writing that
image of \(x\in\Lambda\) as \(\sum_i a_i\iota(\omega_i)\) gives
\[
     \beta(x,\operatorname{Fr} y)
          =\sum_i a_i\beta(\omega_i,y)^p\in R
                    \qquad (y\in\Lambda).
\]
Here \(a_i\in R\), \(\omega_i\in\Omega_0\), and integrality of the
pairings on the right follows from Hodge adaptation. Thus the
projection of \(\Lambda\) to \(H/W\) lies in the full dual of
\(N\).

Let \(P_W\) be the image of \(\Lambda\) in \(H/W\), and
use the pairing to view it in \(N^\vee\). The exact sequence
\(0\to N\to\Lambda\to P_W\to0\) gives, in a symplectic basis,
\[
 v(\det\Lambda)=v(\det N)+v(\det P_W)
       =\operatorname{length}_R(N^\vee/P_W).
\]
Indeed
\[
 v(\det N^\vee)=-v(\det N)
\]
in the dual basis.
Computing the same determinant with the Lagrangian \(F_0\)
instead gives zero, since Hodge adaptation identifies the
projection with \(\Omega_0^\vee\). The resulting inclusion has zero colength, so
\[
 \operatorname{length}(N^\vee/P_W)=0\quad\Longrightarrow\quad P_W=N^\vee.
\]
This is the second adaptation identity
at \((W,N)\).

The comparisons use the Weil pairing and its extension on the
torsion tower, with functoriality for the paired endomorphisms.
The two adaptation identities give the sublattice and full
quotient in each sequence. To descend a
completed lattice, approximate a basis by rational vectors to
precision strictly finer than its elementary divisors. The change
of basis then belongs to the integral general linear group of
the completion. The rational lattice with this completed span
has the same intersections and quotients by faithful flatness.
It modifies the bundle at just the finite set of missing points.
The two sequences extend there by the prescribed lattices and
the split integral action gives their eigensequences.
\label{cor:flag-wall-positivity-7}
Writing \(d_\sigma=\deg D_\sigma\) gives
\[
 d_{\varphi\sigma}=p d_\sigma.
\]
For an orbit of length \(h\), iterating gives
\[
 d_\sigma=d_{\varphi^h\sigma}=p^h d_\sigma .
\]
Since \(p^h>1\), \(d_\sigma=0\), hence
\(\deg U_\sigma=\deg U_{\bar\sigma}\).

\proofstep{Uniform control by the largest subsheaf degree.}
Let \(P\) be the largest degree of a coherent subsheaf of \(U\).
The degrees are integers. For an upper bound, twist the
dual sufficiently to generate it, embed \(U\) into a sum of fixed
ample lines, and take exterior powers. The zero subsheaf gives
\(P\ge0\). Saturate a subsheaf attaining \(P\). Its degree remains
\(P\). Its Frobenius pullback of degree \(pP\) injects in the full
de Rham bundle \(D\), an extension of \(U\) by \(U^\vee\).
If \(P'\subset U\) realizes \(P\), let \(B\) and \(A\) be
the image and kernel of \(F^*P'\to D\to U\), respectively.
Then \[
 0\to A\to F^*P'\to B\to0,\qquad \deg B\le P.
\]
Saturate \(A\subset U^\vee\). The dual of its locally free
quotient embeds in \(U\), so
\[
 \deg A\le\deg U^\vee-\deg(U^\vee/A^{\rm sat})
                  \le-\deg U+P .
\]

Now \(pP=\deg A+\deg B\le-\deg U+2P\), giving
\[
        (p-2)P\le e,\qquad e=-\deg U .
\]
For the root degrees \(r_{\tau i}\) on the resulting curve, set
\(S_{\tau k}=\sum_{i<k}r_{\tau i}\).
The full sum is the difference of the two conjugate Lie
degrees, so \(S_{\tau N}=0\). The middle sum is one full Lie
eigendegree, hence \(-2\sum_\tau S_{\tau j}=e\). Before the
middle, a prefix is the degree of a flag subsheaf and is at
most \(P\). After the middle, reverse the conjugate flag and
use the vanishing full sum to obtain the same upper bound.
The argument gives these properties for a common positive
multiple of the original readings. Rescale \(P,e\) by the
same factor on descent. \label{cor:flag-wall-positivity-8}

Put \(S=\|w\|_\infty\) and suppose
\[
 4f S\le(B-2)\delta_{\tau j}
\]
for each \(\tau\).
Let \(b_p=4fS/(p-2)\), so that \(\delta_{\tau j}\ge b_p\).
Summation by parts and \(P-S_{\tau k}\ge0\) now give
\[
 \begin{split}
 -\sum_{\tau,i} w_{\tau i}r_{\tau i}
 &=\sum_{\tau,k<N,\ k\ge1}\delta_{\tau k}(P-S_{\tau k})
              -P\sum_\tau(w_{\tau0}-w_{\tau,N-1})\\
 &\ge b_p\bigl(|\mathcal T|P+e/2\bigr)-2fSP\ \ge 0 .
 \end{split}
\]
This proves nonnegative degree at every remaining prime.
After adding \(sA\), clearing the line denominators and descending
from the common model we obtain fibrewise nefness of \(M_B(v,s)\)
on each marked component. Together with the archimedean
semipositivity and \eqref{eq:flag-actual-ht} this supplies the
point and positivity estimates for the arithmetic intersection
on both selected covers. \label{cor:flag-wall-positivity-9}
\end{proof}


\subsection{The arithmetic degree near a wall}

The height bounds and vertical positivity can now be applied to weights approaching a chamber wall. The geometric degree tends to zero at the wall. Upper and lower height bounds then control the arithmetic intersection on the wall patch.

\paragraph{The polynomial on the same component.}
Let \(d=d_F\), and write \(\theta(w)\) for the smooth reference
curvature of the root combination. For a weight \(w\) and real
parameter \(s\) define \(I_R(w,s)\) by taking the arithmetic
\((d+1)\)-st power of the root line at \(Q_B\) with its original
truncated metric at \(R\), plus \(sA\), on the component \(c\).
Take the section and field degree normalizations used for
\(\Phi_B\). Since the logarithms of the determinant norms at fixed \(R\)
are linear in the weight, multilinearity gives the degree bound
\[
                      \deg_{w,s} I_R(w,s)\le d+1 .
\] The
cutoff \(R\) in the archimedean metric is independent of the
finite cutoff \(B\).

To bound this polynomial, compare its metric with the perturbation already controlled on horizontal and vertical curves. In a toroidal
chart let \(r\) again denote the logarithmic boundary radius.
The truncated matrix on the first block is the interpolation
\[
 (1-a_R)H_0+a_R H
\]
of the smooth and Hodge Gram matrices.
The same formula on the other block uses the reversed dual
matrices. For \(r\le R^2\) one has
\[
 a_R=(1+(r/R)^\kappa)^{-1}.
\]
The interpolation agrees with
the smooth reference at \(r\ge R^3\). Choose the exponent as
in the chart, \(\kappa>2\mu+2\), where each original form lies
between \(r^{-\mu}\) and \(r^\mu\) times the other. For
\(r\le\sqrt R\), the relative error in the interpolation is at
most \(R^{-1}\), since
\[
 (1-a_R)r^\mu\le R^{(\mu-\kappa)/2}.
\]

The log of each graded norm differs by \(O(1/R)\), by comparing
successive principal determinants. Everywhere its change from
the smooth reference is \(O(\log r)\). Both estimates are
uniform for bounded weights. \label{cor:flag-wall-polynomial-1}

\begin{proposition}[Componentwise arithmetic comparison at the weight wall]
\label{prop:arithmetic-wall-polynomial}
On the fixed component \(c\), the arithmetic power of the
perturbed metric is approximated by the limiting weight polynomial
from \(I_R\). Its logarithmic boundary mass satisfies
\eqref{eq:flag-wall-tail}. The difference of the finite-cutoff
powers satisfies \eqref{eq:flag-wall-finite-cutoff}. Writing \(g_c(v)\) for the geometric degree on the same component, the component measures enter as
\[
 g_c(v)=w_cG(v),\qquad w_c>0,\qquad \sum_c w_c=V.
\]

\end{proposition}

\begin{proof}
\proofstep{Comparing with the intersection in the perturbed metric.}
For the regularized potentials of Section~\ref{subsec:flag-wall-positivity}
write \(\ell=s^{-2}\). On \(r\le\ell\) their distance from the
Hodge potentials is \(O(s)\). Indeed on a ball of radius
\(c e^{-r/2}\) in a crossing chart the new boundary radius is
at most \(r+C\). Differentiating the Hodge Gram ratios bounds
the gradient there by
\[
 C(r+C)^\nu e^{r/2}.
\]

In a logarithmic
coordinate the derivative of a period brings down a factor
\(1/z_i\), bounded by \(Ce^{r/2}\), and the Gram matrices,
their inverses in the eigenframes, and the smooth changes of
frame have polynomial bounds in \(r\). The smoothing radius
\(e^{-\ell}\) times this gradient is \(O(e^{-\ell/3})\).
The bounded regularization corrections of size
\(O((1+\log\ell)/\ell)\) there are also \(O(s)\).
For \(r\ge\ell\) the upper and lower value bounds give a
difference \(O(\log r)\). \label{cor:flag-wall-polynomial-2}

\proofstep{Controlling the logarithmic boundary growth.}
Here is the boundary estimate that permits use of these comparisons
under the intersection integral. A positive curvature of the form
\(\theta+dd^c u\) off the divisor, with
\[
 -C\log r\le u\le0
\]
and \(-C\Omega\le\theta\le C\Omega\),
will be used with \(u\) smooth or tending to \(-\infty\) at the
boundary. Each curvature in the interpolation can be written
as a bounded combination of such forms. More explicitly, for
the truncated Gram metric one adds a fixed multiple of a positive
Poincar\'e form, then subtracts it.

In crossing coordinates the
latter dominates the smooth directions and the forms
\[
 \frac{i\,dz_a\wedge d\bar z_a}
 {|z_a|^2(1+|\log |z_a||)^2}
\]
in the boundary directions. Its potential is a sum of negative
logarithms of the branch radii, comparable to \(-\log r\)
up to fixed factors. The lower curvature bound for the interpolated
metric uses the blockwise decreasing weights. In the chart curvature
formula the two end metrics have bounded normalized curvature
and their connection difference is bounded in the Poincar\'e
metric. Inserting the interpolated Gram matrix in
\(\bar\partial(H^{-1}\partial H)\), the terms with a derivative
of \(a_R\) are controlled by the logarithmic derivatives of the
cutoff.

On \(r\le R^2\) these are obtained from
\[
 da_R=-\kappa a_R(1-a_R)\,dr/r.
\]
On the terminal transition the
squared smooth cutoff and \(\kappa>2\mu+2\) bound the terms
involving the inverse Gram matrix. The flag restriction adds
the nonnegative second fundamental form for each proper step.
This is the interpolation curvature estimate on the adapted
toroidal charts. The positive correction can be enlarged so its
potential dominates the absolute log growth. For the \(s\)-metric
use its semipositivity and uniform upper bound. Majorization of the
Hodge potential gives the same logarithmic lower bound off the
divisor, uniformly in \(s\). Constant changes in potentials leave
their curvatures unchanged. \label{cor:flag-wall-polynomial-3}

For a product \(d\mu\) of \(d\) of the indicated positive forms,
the mass is bounded, and for \(D\ge e\)
\begin{equation}\label{eq:flag-wall-tail}
 \int_{\{r\ge D\}}\log r\,d\mu
       \ \ll\ \frac{(1+\log D)^{d+1}}{D}.
\end{equation}
To verify the bound on a shell \(D\le r\le 2D\), truncate
each \(u\) from below by a smooth convex maximum with
\(-C\log(2D)-2\), agreeing with \(u\) on the shell.

Adding
\(C\Omega\) makes the truncated curvatures positive everywhere.
They are smooth through the divisor, since a potential tending
to \(-\infty\) is then constant near it. Their mixed masses
against powers of \(\Omega\) are uniformly bounded by Stokes'
theorem applied to the fixed smooth parts. Represent each
truncated potential up to a constant by
\(0\le v\le C'(1+\log D)\). Take a smooth increasing concave
truncation \(r_D\) of \(r\), equal to \(r\) up to \(2D\) and
constant beyond \(3D\), with derivative at most one.
The logarithmic divisor formula gives \(dd^c r_D\le C\Omega\).
Also
\[
 \int r_D\Omega^d\le\int r\Omega^d=O(1).
\]

Write one truncated curvature as \(dd^c v+\alpha\) where
\(\alpha\) is smooth and uniformly bounded above by \(C\Omega\).
For a closed mixed positive product \(P\) of complementary
degree, Stokes gives
\[
 \int r_D(dd^c v+\alpha)\wedge P
   =\int v\,dd^c r_D\wedge P+\int r_D\alpha\wedge P
   \le C'(1+\log D)\int\Omega\wedge P+
                              C\int r_D\Omega\wedge P.
\]
The first integral on the right is a uniformly bounded mass.
The second replaces one truncated curvature by \(\Omega\).
Iteration terminates at \(\int r_D\Omega^d=O(1)\), so for a mixed
product \(P_d\) of \(d\) such curvatures the shell bound is
\[
 \int_{\{D\le r<2D\}}P_d
  \ \le\ D^{-1}\int r_D P_d
  \ \ll\ (1+\log D)/D \qquad (d\ge1).
\]
Summing with weight \(\log(2D)\) on dyadic shells gives
\eqref{eq:flag-wall-tail}.
The mass on the part of bounded radius follows by the same
truncation and ordinary mixed-mass bound. This applies also
to absolute values of integrals of the bounded combinations,
by expanding their products. \label{cor:flag-wall-polynomial-4}

\proofstep{The difference of arithmetic powers.}
Changing just the archimedean metric in the arithmetic power
telescopes into \(d+1\) integrals of the potential difference
against mixed curvatures, with the common component
normalization. Write \(\varphi_{R'},\varphi_R\) for the normalized root
potentials in the two readings and \(\mathcal M_i\) for the products
of \(i\) curvatures from \(R'\) and \(d-i\) from \(R\).
Use the \(O(1/R)\) bound on \(r\le\sqrt R\) and the weighted
tail \eqref{eq:flag-wall-tail} on its complement:
\[
 \sum_{i=0}^d\left|\int (\varphi_{R'}-\varphi_R)\,\mathcal M_i\right|
   \ll R^{-1}+\frac{(1+\log R)^{d+1}}{\sqrt R}.
\]
The same component normalization is used in the two degrees,
and the bound is uniform in \(R'\ge R\), \(B\), the bounded
dominant weights and \(0\le s\le1\). Thus they converge on an
open set of strictly decreasing weights and \(0<s<1\).
Evaluation on that open set separates polynomials of bounded
degree.

A finite subset of evaluations forms a basis of their
dual, so that all the coefficients converge. Call the limiting
polynomial \(P_{B,c}\). At \(s=0\) its component sum is
\(\Phi_B(L_w^{d+1})\) by the defining limit, first on the open
chamber and then identically. Compare the \(s\)-metric to the
truncated metric at \(R\ge\ell^2\). Their difference is \(O(s)\)
on \(r\le\ell\) and \(O(\log r)\) elsewhere. Putting \(\ell=s^{-2}\) in the tail estimate gives
\[
 \ell^{-1}(1+\log\ell)^{d+1}
 =s^2(1+2\log(1/s))^{d+1}=O(s)
\]
as \(s\downarrow0\), for the fixed dimension \(d\). Taking \(R\to\infty\)
proves that the actual component intersection for \(M_B(v,s)\)
differs from
\[
 P_{B,c}(v-sd_0,s)
\]
by \(O(s)\), uniform in
\(B\) on the fixed model. \label{cor:flag-wall-polynomial-5}

\paragraph{Dependence on the finite cutoff.}

For wall estimates obtained by varying the cutoff to determine the functional at a fixed cutoff, the difference of the two arithmetic degrees must be tracked prime by prime.

Compare \(B\) with \(2\) on a common refinement and divide by
the corresponding generic cover degrees. The changed primes
are good and their root log norms differ by at most
\(\log p/(p-1)\) in the normalization at \(p\). For a combination
\(x\) at \(B\) and \(y\) at \(2\) with the same bounded weights
and ample summand, the generic and archimedean lines coincide.
The vertical log bound for \(x-y\) is at most the sum of the
absolute weight coefficients times the root bound. On a list
of \(d\) lines nef over the changed primes its mixed intersection
is therefore bounded in absolute value by
\[
 C\sum_{p\le B}\log p/(p-1)
\]
times their generic mixed degree.
Indeed the local vertical divisor is between the two corresponding
multiples of the whole fibre. Restrict the nef lines to the
components and sum their nonnegative mixed degrees with these
multiplicities.

For clarity, \(y\) itself need not be nef at the newly changed
primes. On the model at \(2\) take a fixed relatively ample
\(H\) large enough that \(H\pm x_i\) are relatively nef for
every reference root \(x_i\). For \(\|w\|_\infty\le1\) and
\(s\ge0\), \(y+h\) is in the nef cone, where \(h\) is the
number of roots times \(H\). This follows by expanding each
\(H-w_i x_i\) with coefficients \((1\pm w_i)/2\).

At the changed primes \(x\) is nef by the Hodge--Tate step
comparison. Use
\[
 x^{d+1}-y^{d+1}=\sum_{i=0}^d(x-y)x^i y^{d-i},\qquad
 y^{d-i}=\sum_{j=0}^{d-i}\binom{d-i}{j}(-1)^j
                                 (y+h)^{d-i-j}h^j .
\]
After taking absolute values of these coefficients, every
local mixed degree uses only the nef lines \(x,y+h,h\).
In the bound by generic degrees replace \(x\) by \(y\),
whose generic class agrees. The sum over \(j\) is
\(\deg_{\mathrm{geom}}(y^i(y+2h)^{d-i})\). It is bounded on the
fixed compact coefficient range independently of the refinement,
after division by the cover degree. This yields
\begin{equation}\label{eq:flag-wall-finite-cutoff}
 \lvert I_R(w,s;B)-I_R(w,s;2)\rvert\ll\log B .
\end{equation}
The same inequality holds in the limit. Applying the same finite matrix of evaluations on the open
chamber now bounds every coefficient of \(P_{B,c}\) by \(O(\log B)\).
\label{cor:flag-wall-polynomial-6}

\paragraph{The geometric degree.}
The canonical root class \(L_v\) is nef for dominant \(v\).
Indeed the regularizations above approximate it by smooth forms
with curvature at least \(-\varepsilon\Omega\), with
\(\varepsilon\downarrow0\). The shifted geometric class
\(M=L_{v-sd_0}+s A\) is nef for dominant \(v-sd_0\). In the
polarization of the toroidal model \(L_{d_0}-A\) has the class
of an effective boundary divisor \(D_0\), with the multiplicities
of its local boundary section. The telescoping formula gives
\[
  \deg L_v^d-\deg M^d
     =s\sum_{i=0}^{d-1} D_0\cdot L_v^i M^{d-1-i}\ge0 .
\]
Each summand is nonnegative by restriction to the effective
components and approximation by positive forms. Polynomial
continuity gives convergence to \(\deg L_v^d\) as \(s\downarrow0\).

For the canonical extension of the homogeneous bundles over a
smooth toroidal compactification of a neat arithmetic quotient,
Mumford's proportionality theorem identifies the Chern evaluation
with the compact-dual evaluation times the component volume in
the proportionality normalization. Passing to the flags uses the
flag-bundle pushforward on both sides, so the weighted root
powers have the same factor. Dividing by the degree of a common
neat cover gives the same formula on the selected component.
It holds first at integral weights and then identically by
multilinearity. The resulting upper degree is \(w_cG(v)\),
where \(w_c>0\) and \(\sum_c w_c=V\).
\cite{mumford-flag-wall}\label{cor:flag-wall-polynomial-7}
\end{proof}


\begin{lemma}[Vanishing on split walls]\label{lem:flag-wall-small}
Fix the marked flag model, and let \(B\) tend to infinity through
integers. At each embedding take the wall patch of weights
\(\|u\|_\infty\le1/2\) with \(u_{j-1}=u_j\) at that embedding
and all other successive gaps at least \(1/(4N)\).
Let \(\mathcal W\) be the union of these patches. For the fixed marked model,
\[
 \sup_{u\in\mathcal W}|H_B(u)|=O\!\left(\frac{\log B}{B}\right).
\]

\end{lemma}

\begin{proof}
\proofstep{Squeezing the arithmetic self-intersection.}
For \(C,a>0\), maximize the error term in \eqref{eq:flag-actual-ht}:
\[
 \sup_{z>0}(C\log z-az)=C\log(C/a)-C,\qquad z_{\max}=C/a .
\]
In the interior substitute \((a,z)=(t/4,2+h)\); on the boundary
substitute \((t^2,2+C+h_A)\). For \(B\ge2\)
and \(0<t\le1/2\), the minimum height is consequently at least
\(-C'\log(B/t^2)\), and the height on the dense set is at most
\(C'\log B\). On the semipositive fibrewise nef line of
Section~\ref{subsec:flag-wall-positivity} these squeeze its
component self-intersection between the two bounds multiplied
by \(n\) times its geometric degree. If \(I\) is the arithmetic self-intersection and \(d_{\rm geom}\)
the geometric degree, shifting every height by \(-e\) gives
\[
                         I_{-e}=I-ne\,d_{\rm geom}.
\]

At the lower bound the shifted line is
arithmetically nef and has nonnegative self-intersection. At an
upper bound increased by any \(\delta>0\) the shifted line has
strictly negative height on a dense set. A nonzero small integral
section of a positive power would have nonnegative height off
its zero locus, contradicting density. Its arithmetic volume is
therefore zero. The arithmetic volume inequality for a smooth
semipositive, relatively nef Hermitian line bounds its
self-intersection above by that volume. Let \(\delta\downarrow0\).
The change to an integral power on an inducing model multiplies both
sides of each comparison by the same positive factor. The constants
here are uniform for the bounded weights and the auxiliary cutoff
on each fixed model. \label{cor:covolume-wall-shape-1}

\proofstep{Perturbing the weight across the wall.}
Set \(c_0=1/(4N)\). At one embedding \(\tau_0\) use the compact wall patch 
\[
 \begin{gathered}u_{\tau_0,j-1}=u_{\tau_0,j},\qquad \|u\|_\infty\le\tfrac12,\\ u_{\tau,i}-u_{\tau,i+1}\ge c_0\qquad((\tau,i)\ne(\tau_0,j-1)).\end{gathered}
\]
 Start with rational \(u\) there.
Take
\[
 t=5f/B
\]
and add \(2t\) to all entries of the first
block at \(\tau_0\) only, obtaining \(v\). The distinguished
gap is exactly \(2t\) and the other gaps stay fixed. For large
\(B\),
\[
 t^2\|d_0\|_\infty\le t/2,
\]
so all adjacent gaps
of \(v-t^2d_0\) are at least \(t\). Its norm is at most one.
Since
\[
 (B-2)t\ge4f
\]
for \(B\ge10\), this gives the nef
inequality. The Vandermonde product gives \(0\le G(v)\ll t\)
because its distinguished difference is \(2t\) and all other
factors are bounded. Also \(\log(B/t^2)\le3\log B\).
The height squeeze on each component now bounds the absolute
total self-intersection by \(O(\log B/B)\).

\proofstep{Returning to the wall.}
Compare the polynomial readings at \((v-t^2d_0,t^2)\) and
\((u,0)\). On the unit cube a monomial of degree \(m\) has
Lipschitz constant at most \(m\) for the sup norm, by telescoping
its factors. There are finitely many monomials and components.
The change costs \(O(t\log B)\), and agreement with the actual
intersections costs \(O(t^2)\). Thus
\begin{equation}\label{eq:flag-wall-small}
   \sup_{\tau_0}\ \sup_{u\text{ on the wall patch at }\tau_0}
                 |H_B(u)|\ll \log B/B .
\end{equation}
For real weights use componentwise rounding down on the mesh
\(1/(4Nm)\), \(m\) integral. The endpoints of the interval and
the minimum gap are multiples of that mesh, so both inequalities
and the wall equality persist. Continuity gives the bound on
the closed patch.

\end{proof}

\begin{established}[Wall vanishing, alternation and permutation of the embeddings]\label{hyp:flag-wall-symmetry}
For the full marked object of Definition~\ref{hyp:covolume-flag-degrees}, the reference polynomial \(H\) is nonnegative in the dominant chamber and vanishes when a split gap closes. For any adjacent exchange \(s\) within a block and each cutoff \(B\ge2\),
\[
       (H_B-VH)(su)=-(H_B-VH)(u).
\]
There is a transitive permutation group of \(\mathcal T\) preserving \(H_B,H,G\) for all \(B\). The arithmetic degree in this assertion sums over the marked component set.
\label{cor:covolume-wall-shape-2}
\end{established}
\begin{proof}
On the reference product of complete flags write an integral dominant line as a tensor product of nonnegative powers of the Pl\"ucker lines and determinant lines. The quotient metrics are induced by the standard Hermitian norm. Their Pl\"ucker coordinates are integral sections, of sup norm at most one, generating each line, so these lines are arithmetically nef. The standard determinant frame has norm one. When two adjacent weights agree the line, including its metric, descends to the partial flag omitting that step. If \(\pi\) is the resulting projection, then
\[
 \widehat c_1(\pi^*\overline L)^{d_F+1}
           =\pi^*\widehat c_1(\overline L)^{d_F+1};
\]
its arithmetic degree is zero by the projection formula since the relative dimension is positive. Rational approximation inside the face and homogeneity give the assertions for real weights. Every permutation of equal flag factors preserves the reference height and the product for \(G\).

For reference the two conclusions just proved read
\[
       H(u)\ge0\ (u\text{ dominant}),\qquad
       u_{\tau,j-1}=u_{\tau,j}\Longrightarrow H(u)=0
       \quad(u\text{ dominant}).
\]

For two adjacent roots write \(t=(x+y)/2\) and \(v=y-x\). The even part under exchange consists of monomials \(v^{2i}M(t,\ldots)\). The rank-two formula in Proposition~\ref{lem:arithmetic-flag-alternation} evaluates each as \(h_i\) times its geometric predecessor on both flag spaces. The latter geometric degrees differ by the volume factor \(V\). If \(\operatorname{Sym}_s=(1+s)/2\), the equality on the even monomials gives
\[
             \operatorname{Sym}_s(H_B-VH)=0 .
\]

For the marked arithmetic functional let \(K\) be the marked imaginary quadratic field and \(\Psi\) the host Galois subgroup fixing it. The selected embeddings form a torsor under right translation by \(\Psi\): for two labels \(\tau,\tau'\) take \(g=\tau^{-1}\tau'\) in torsor coordinates. The two half-level ideals are extended from the marked quadratic order, so \(g\) fixes them. Precomposing the PEL endomorphism action with \(g\), and transporting the labels of the marking by the same map, gives an invertible change of moduli data. Its inverse uses \(g^{-1}\). For the twisted endomorphism action \(\iota^g=\iota\circ g\), its eigenspaces satisfy
\[
   \operatorname{Lie}(A,\iota^g)_\tau
           =\operatorname{Lie}(A,\iota)_{\tau\circ g^{-1}} .
\]
The transport preserves polarization, complex conjugation and the flag on each eigenspace. Write \(C\mapsto gC\) for the induced bijection of components.

Use the transported toroidal fan for the conjugate moduli problem. At a finite place transport also the test lattice: the cutoff selects rational primes and the de Rham or Hodge--Tate construction commutes with the transported action. At infinity the polarization carries the flag metrics and their cutoffs isometrically. Thus for any fan \(\Sigma\) and smoothing radius \(R\), change of variables gives
\[
 \sum_C\widehat{\deg}_{C,\Sigma,B,R}(L_{u\circ g}^{n})
       =\sum_C\widehat{\deg}_{gC,g\Sigma,B,R}(L_u^n).
\]
It suffices to use a common refinement of the fans. On that refinement the geometric and finite intersections satisfy the projection formula. In the archimedean difference of the two smoothed powers, each term contains one difference of cutoff potentials, supported toward the boundary; the integrable bounds and telescoping estimate in the proof of Proposition~\ref{prop:arithmetic-wall-polynomial} make its integral tend to zero at fixed \(B\). Taking \(R\to\infty\) gives \(H_B(u\circ g)=H_B(u)\), as required.
\end{proof}

\begin{construction}
\[
\begin{array}{ll}
\text{within-block alternation}&\prod_{i<k\text{ in a common block}}(u_i-u_k)\\
\text{vanishing on the cross-block walls}&\prod_{i<j\le k}(u_i-u_k)\\ \hline
\text{remaining block sums}&\alpha_\tau\!\sum_{i<j}u_{\tau i}
                      +\beta_\tau\!\sum_{i\ge j}u_{\tau i}
\end{array}
\]
\caption{How the vanishing conditions factor a wall polynomial. At each embedding the last quotient is linear and symmetric within each of the two blocks.}\label{fig:wall-kernel-steps}
\end{construction}

\begin{proposition}[The two coefficients left by a split wall]\label{prop:flag-two-moving-shifts}
For the fixed marked model of Lemma~\ref{hyp:flag-wall-symmetry}
there are real \(a_B,b_B\) such that
for each fixed homogeneous polynomial \(F\) of degree \(n\),
\begin{equation}\label{eq:flag-moving-shifts}
 \Phi_B(F)=V\bigl(h_{\rm fl}(F)
                  -a_B\mathsf S(F)-b_B\mathsf T(F)\bigr)
                     +O_F(\log B/B).
\end{equation}
The implied constant may depend on the fixed model and \(F\).
The coefficients \(a_B,b_B\) are common to all \(F\).
\end{proposition}

\begin{proof}
\proofstep{Factoring by the wall equations.}
In the finite-dimensional space of block-antisymmetric forms
of degree \(n\), take the common kernel of the evaluations on
all wall patches. Its elements are precisely
\[
  G(u)\sum_\tau\left(\alpha_\tau\sum_{i<j}u_{\tau i}
                    +\beta_\tau\sum_{i\ge j}u_{\tau i}\right).
\]
Indeed, antisymmetry supplies each within-block difference as
a factor. A patch contains a relatively open subset of its
hyperplane. To see this, take all other gaps strictly between \(1/(4N)\)
and \(1/(2N)\) and center the entries. A polynomial vanishing
there is divisible by the distinguished cross-block difference.
Permuting either block gives every cross-block difference.
The distinct linear factors are relatively prime in pairs
and their product has degree \(d_F\). The linear quotient
must have constant coefficients within each block because
both the dividend and the Vandermonde are alternating there.

Let \(\mathcal K\) be the common kernel of the wall evaluations
calculated above. Fix a complement \(\mathcal L\), and choose
wall evaluations that form a basis of \(\mathcal L^\vee\).
For \(Q_B=H_B-VH\), decompose
\[
 Q_B=Q_{B,\mathcal K}+E_B
\]

with \(E_B\in\mathcal L\). In a fixed basis \(f_1,\ldots,f_h\) of \(\mathcal L\), let \(A_{ij}=f_j(u_i)\) for the chosen wall points \(u_i\). The evaluation functionals span \(\mathcal L^\vee\): a vector they all annihilate belongs to \(\mathcal L\cap\mathcal K=0\). We may therefore take \(A\) invertible. Writing \(c_B\) for the coefficient vector of \(E_B\) gives
\[
 A c_B=(Q_B(u_i))_{i=1}^h,\qquad
 \|c_B\|_\infty\le\|A^{-1}\|_\infty\max_i|Q_B(u_i)|.
\]
At this fixed model the matrix is independent of \(B\). Equation~\eqref{eq:flag-wall-small} thus gives
\[
 \|E_B\|_{\rm coeff}=O(\log B/B).
\]

Average both summands over the finite transitive relabeling
group. This leaves their sum equal to \(Q_B\), preserves
\(\mathcal K\), and keeps the coefficient bound because each
relabeling permutes the monomials. In the averaged kernel part
the first-block coefficient is the same at every embedding,
and the second-block coefficient is likewise constant.
Write them as \(nV(a_B+b_B)\) and \(nVb_B\), respectively.
The derivative degrees in Section~\ref{sec:covolume-flag-test}
identify these two terms with the \(\mathsf S,\mathsf T\) terms
in \eqref{eq:flag-moving-shifts}.

\proofstep{Passing from weights to test polynomials.}
To see this first on \(L_u^n\), its derivative in \(x_{\tau i}\)
has degree reading
\[
 -nG(u)u_{\tau i}.
\]
For a monomial
\(x^m\) of degree \(n\) the error in its functional value is
the coefficient of \(u^m\) in the polynomial error divided by
\((-1)^n\binom{n}{m}\). Each is \(O(\log B/B)\). Linearity
gives \eqref{eq:flag-moving-shifts}. No bound for the size of the
inverse evaluation matrix as \(N\) varies is needed in this limit.

\end{proof}

\paragraph{Passing to a fixed real cutoff.}
Write \(P_{\tau,k}=P^-_{\tau,k}+P^+_{\tau,k}\).
Every homogeneous \(F\) of degree \(n\) has a presentation
\begin{equation}\label{eq:flag-power-sum-presentation}
 F=\sum_\tau\sum_{k=1}^N P_{\tau,k}Q_{\tau,k},
 \qquad \deg Q_{\tau,k}=n-k .
\end{equation}
with homogeneous cofactors (possibly zero). For a single
embedding the ideal of the first \(N\) power sums equals the
ideal of positive elementary symmetric functions by Newton's
identities over \(\R\). Modulo this ideal the complete symmetric
function of degree \(i+1\) in \(x_i,\ldots,x_{N-1}\) vanishes because
its generating series there is \(\prod_{a<i}(1-x_a z)\).
Thus one reduces \(x_i^{i+1}\) to smaller powers of \(x_i\)
with later variables. In increasing order of \(i\) this leaves
monomials of degree at most \(\binom N2\) per embedding. Taking
the homogeneous part of the reductions gives the presentation.
\label{cor:covolume-wall-shape-3}

\begin{proposition}[The fixed-cutoff expansion]\label{prop:flag-fixed-prime-expansion}
Fix \(B_0\) and a marked model satisfying the prime-change
formulas of Section~\ref{sec:covolume-flag-test}.
For a homogeneous \(F\) of degree \(n\), choose any power-sum
presentation \eqref{eq:flag-power-sum-presentation}. There are constants \(a,b_0\), independent of \(F\),
such that
\[
 \Phi_{B_0}(F)=V\bigl(h_{\rm fl}(F)
                 -a\mathsf S(F)-b_0\mathsf T(F)\bigr)
                   +\sum_{p\notin\mathcal P_{B_0}}R_p(Q).
\]
For each fixed presentation the series is absolutely convergent.
\end{proposition}

\begin{proof}
\proofstep{Corrected coefficients at the fixed real cutoff.}
Put \(\mathcal P'=\mathcal P_B\setminus\mathcal P_{B_0}\) for
\(B\ge B_0\). At a new good prime the de Rham minus Hodge--Tate
comparison of the power sum has a reference term
\(-k C(P^-_{\tau,k-1}Q_{\tau,k})\log p/(p-1)\).
Here \(C=V\deg_{\rm fl}\), including the exponent-zero convention.
For an embedding label \(\alpha\), write \(\mathrm{Lie}_\alpha\) for the corresponding Lie eigenspace bundle and \(\operatorname{ch}_{k-1}(\mathrm{Lie}_\alpha)\) for its signed power sum divided by \((k-1)!\). Let \(\varphi_p\) be the Frobenius permutation of the labels. The successive Frobenius quotients contribute
\[
 k!\log p\sum_{l\ge1}p^{-kl}
 C\big(\operatorname{ch}_{k-1}(\mathrm{Lie}_{\varphi_p^l\tau})
                                      Q_{\tau,k}\big).
\]
These formulas follow from the periodic quotient
identity in Section~\ref{sec:covolume-flag-test} by expanding the
geometric series around the Frobenius orbit.

Changing the
cofactor's flags after the power sum has changed contributes zero
by vertical restriction of the trivialized Tate bundle.
The reference terms sum to \(-V\mathsf S(F)\log p/(p-1)\) because
terms with a derivative on the cofactor still have a full power
sum as factor and have zero geometric degree. For \(k=1\) the
Lie rank is \(j\) or \(N-j\), according to the Frobenius type.
Let
\[
 \kappa_p=\sum_{l\ge1}p^{-l}\operatorname{rk}(\mathrm{Lie}_{\varphi_p^l\tau})
\]
for a plus embedding \(\tau\).
This is the same for every such \(\tau\). They are indexed by
the host subgroup of the Galois group, and multiplying by any
of its elements preserves membership in the subgroup and its
conjugate coset.

Since the total derivative of \(P_{\tau,1}\)
at \(\tau\) is \(N\), the \(k=1\) terms sum to
\(V\mathsf T(F)\kappa_p\log p/N\). Denote the remaining sum at
\(p\), over \(\tau\) and \(k\ge2\), by \(R_p(Q)\).
Telescoping over the new primes and taking the smoothing limit
for this finite set gives the exact change from \(B\) to \(B_0\).
For fixed \(Q\) its higher-order terms satisfy
\(\abs{R_p(Q)}\ll_Q \log p/p^2\) since the possible Chern
readings are finite and
\[
 \sum_{l\ge1}p^{-kl}=1/(p^k-1).
\]

\label{cor:covolume-wall-shape-4}

The finite telescoping identity and
Proposition~\ref{prop:flag-two-moving-shifts} combine to give
\[
 \Phi_{B_0}(F)=V(h_{\rm fl}(F)-A_B\mathsf S(F)-D_B\mathsf T(F))
                    +\sum_{p\in\mathcal P'}R_p(Q)+o_F(1),
\]
where the reference and rank-one prime terms have been incorporated
into the corrections
\[
 A_B=a_B+\sum_{p\in\mathcal P'}\frac{\log p}{p-1},
 \qquad D_B=b_B-\frac1N\sum_{p\in\mathcal P'}\kappa_p\log p .
\]
The bound \(R_p(Q)=O_Q(\log p/p^2)\) applies to each fixed
presentation, so it remains to take the limits of \(A_B,D_B\).
For the two calibration polynomials the derivative rows are
\((2I,0)\) and \((j,N)\). In these two rows the preceding
identity says that both \(2I A_B\) and \(jA_B+ND_B\)
converge, because every other term has a limit. Since
\(I\ne0\) and \(N>0\), this first gives \(A_B\to a\), then
\(D_B\to b_0\). Substituting for an arbitrary \(F\) yields
\begin{equation}\label{eq:flag-fixed-shifts}
 \Phi_{B_0}(F)=V(h_{\rm fl}(F)-a\mathsf S(F)-b_0\mathsf T(F))
                        +\sum_{p\notin\mathcal P_{B_0}}R_p(Q).
\end{equation}
The constants came from the two calibration rows. The finite
telescoping identity holds for all power-sum presentations, so
the limiting formula does also.

\end{proof}

\begin{paramdefinition}[Weight gaps and cutoff for the quantitative prime tail]\label{hyp:flag-prime-tail}
For the numerical application suppose \(u\) is strictly
dominant and \(g>0\) bounds below each split gap and the
distance from each entry in the smaller block to every other
entry at the same embedding. In a tie take the first block.
Set \(j_s=\min(j,N-j)\), \(U_\tau=\max_i u_{\tau i}-\min_i
u_{\tau i}\), and take
\[
 B_0=4fN^2(1+\|u\|_\infty/g),\qquad F=L_u^n .
\]
Since \(f(N-1)>2\), this cutoff is at least \(48\).
Qualitative convergence for a fixed presentation in
\eqref{eq:flag-fixed-shifts} must now be strengthened. The
individual polynomials \(Q_{\tau,k}\) need no size bound.
For \(2\le k\le N\), consider the normalized intersection numbers
\[
 K(\tau,k;\sigma,i)=
  \frac{k\deg_{\rm fl}(x_{\sigma i}^{k-1}Q_{\tau,k})}{nG(u)} .
\]
\end{paramdefinition}

\begin{proposition}[Uniform absolute bound for the prime tail]\label{prop:flag-uniform-prime-tail}
Let \(u\) satisfy Definition~\ref{hyp:flag-prime-tail}. Take
\(B_0=4fN^2(1+\|u\|_\infty/g)\), \(j_s=\min(j,N-j)\),
and \(U_\tau=\max_i u_{\tau i}-\min_i u_{\tau i}\).
For \(F=L_u^n\), the tail in \eqref{eq:flag-fixed-shifts}
satisfies
\[
 \frac{\sum_{p\notin\mathcal P_{B_0}}|R_p(Q)|}{nVG(u)}
                         \le 8\log4\, j_s\sum_\tau U_\tau .
\]
The bound is independent of the choice of homogeneous cofactors.
\end{proposition}

\begin{proof}
We give the coefficient calculation underlying the uniform bound.
For \(\rho=(N-1,N-2,\ldots,0)\) and a decreasing exponent
list \(b\) let \(D_b(z)=\det(z_i^{b_l}/b_l!)_{i,l}\). Put
\[
   D=D_\rho,\qquad
   J_k=\sum_{h=0}^{k-1}(-1)^hD_{\rho+(k-h,1^h)} .
\]
Partitions in this formula are padded to length \(N\).
The alternating coefficient of the power-sum presentation gives
\[
 K(\tau,k;\sigma,i)=
 \begin{cases}
 -\partial_i^{k-1}J_k(u_\tau)/D(u_\tau),&\sigma=\tau,\\
 -\big(J_k(u_\tau)/D(u_\tau)\big)
    \partial_i^{k-1}D(u_\sigma)/D(u_\sigma),&\sigma\ne\tau .
 \end{cases}
\]
Here is one derivation. Multiply the presentation by
\(x_{\sigma i}^{k-1}\), alternate at every embedding and
extract the coefficient at the exponent list \(\rho\) at the
embeddings other than \(\tau\).

At \(\tau\) use the signed sum
of the hook lists \(\rho+(k-h,1^h)\) divided by \(k\).
Use the same overall sign as in \(\deg_{\rm fl}\).
After factoring off each Vandermonde these are ordinary Schur
coefficients. The hook sum pairs with the power sum \(p_k\)
in the scalar product for which Schur functions are orthonormal, since
\[
 p_k=\sum_h(-1)^h s_{(k-h,1^h)}.
\]

Multiplication by \(p_l\) is adjoint to \(l\partial/\partial p_l\),
by the power-sum scalar product. Thus only the factor
\(P_{\tau,k}\) survives and its multiplier \(k\) cancels the
division. These identities apply in symmetric degree at most
\(k\le N\) at the selected embedding. At any other embedding
the coefficient sought has symmetric degree zero.
Expanding the other side by the multinomial theorem gives a
determinant \(D\) at each ordinary embedding and \(J_k\) at
\(\tau\). Inserting \(x_{\sigma i}^{k-1}\) lowers its exponent
in the multinomial coefficient, which differentiates row \(i\)
there \(k-1\) times. Division by \(nG(u)\) cancels the common
factorials and the unmodified determinants. The sign comes from
\[
 L_u=-\sum u_{\tau i}x_{\tau i}.
\]
\cite{macdonald-symmetric}\label{cor:covolume-wall-shape-5}

\proofstep{Translation invariance of the normalized derivative.}
The quotient \(J_k/D\) is homogeneous of degree \(k\)
and its coefficient absolute sum is at most \(2^{k-1}/k!\).
In fact the \(h\)-th determinant ratio before the column
factorials is the hook Schur polynomial, with nonnegative
coefficients and value \(\binom{N+k-h-1}{k}\binom{k-1}{h}\)
at \((1,\ldots,1)\), by the hook-content formula. Multiplication by the column factorial ratio and summation over the
hooks give the coefficient bound (a missing hook contributes zero)
\[
 \sum_h\frac{(N-1-h)!}{(N+k-h-1)!}
       \binom{N+k-h-1}{k}\binom{k-1}{h}
 \le\frac1{k!}\sum_{h=0}^{k-1}\binom{k-1}{h}=\frac{2^{k-1}}{k!},
\]
where the sum on the left is over the hooks occurring in the determinant.
Both \(D\) and \(J_k\) are invariant under common translation
for \(k\ge2\).

For \(D\) this is the Vandermonde formula.
For \(J_k\), differentiating in the translation lowers one
column exponent by one with coefficient one. Terms with a
repeated exponent vanish. The surviving ways of removing
a box from the arm or leg give the same determinant for
successive hooks with opposite signs. This also proves the
translation identity for every partial derivative.

Translate the \(\tau\)-variables by \(-u_{\tau i}\), so their
absolute values are at most \(U_\tau\).
Translation invariance and the coefficient bound give the first
estimate below; the separation of the smaller-block coordinate \(i\)
gives the second:
\[
 \left|\partial^l(J_k/D)\right|
       \le \frac{2^{k-1}}{k!}k^l U_\tau^{k-l},\qquad
 \left|\frac{\partial_i^a D(u_\sigma)}{D(u_\sigma)}\right|\le(N/g)^a .
\] Indeed,
differentiate the product of its \(N-1\) factors involving
that index and use the prescribed separation at \(\sigma\).
Leibniz' formula in the equal-embedding case and the product
formula at distinct embeddings now give
\[
 |K(\tau,k;\sigma,i)|\le
 \frac{2^{k-1}}{k!}U_\tau\,[N(1+U_\tau/g)]^{k-1}
 \le U_\tau(B_0/2)^{k-1}.
\]
This use of translation is essential for the spread in the
last estimate. A common shift of the weight at an embedding
can make its sup norm much larger than its spread.

The \((k,l,\tau)\)-summand of \(R_p(Q)/(nVG(u))\) is
\(p^{-kl}\log p\) times a signed sum of the normalized
coefficients \(K(\tau,k;\sigma,i)\) in the Frobenius block.
The full power sum of degree \(k-1>0\) has zero pairing
in the flag ring, so replacing this block by its complement
only reverses the sum's sign. Thus there are at most \(j_s\)
entries to count at each \((k,l,\tau)\). Using the preceding
coefficient bound, first sum in \(l\), then in \(k\):
\[
 \begin{split}
 \frac{|R_p(Q)|}{nVG(u)}
 &\le j_s\sum_\tau U_\tau\,\log p
     \sum_{k=2}^{N}(B_0/2)^{k-1}\sum_{l\ge1}p^{-kl}\\
 &\le 2 B_0 j_s\sum_\tau U_\tau \frac{\log p}{p^2}
                                      \qquad(p>B_0).
 \end{split}
\]

For the second inequality use
\(\sum_{l\ge1}p^{-kl}\le2p^{-k}\) and
\(\sum_{k\ge2}(B_0/(2p))^{k-2}\le2\).
All primes at most \(B_0\) are in \(\mathcal P_{B_0}\).
On \(2^a B_0<p\le 2^{a+1}B_0\), the Chebyshev estimate
\(\sum_{p\le t}\log p\le t\log4\) bounds the sum of
\(\log p/p^2\) by \(2^{1-a}\log4/B_0\).
Sum over \(a\ge0\) to obtain \(4\log4/B_0\). Multiplication
gives the claimed tail bound.

\end{proof}

\paragraph{The component mean.}

The polynomial estimates concern the marked components separately. The slope estimate on the cover uses their mean with equal weights.

\begin{proposition}[Lower bound for the mean of the component slopes]\label{prop:flag-component-mean-lower}
Use \(B_0,u\) of Definition~\ref{hyp:flag-prime-tail} on the
marked component models. Let \(s_c\) be the limiting normalized
slope computed from the component degree in Proposition~\ref{prop:arithmetic-wall-polynomial}, and put
\[
 \ell_{B_0}=\log(B_0\operatorname{rd}_E)
                      +|\RePart (L'/L)(1,\chi_K)|.
\]
For rational \(u\) in the strict chamber one has
\[
 \frac{1}{\#\{c\}}\sum_c s_c
 \ge-C_*\left(\ell_{B_0}\sum_{\tau,i}|u_{\tau i}|
                              + j_s\sum_\tau U_\tau\right).
\]
The slope average uses equal component weights; the sum of
the common component weights is \(V\).
\end{proposition}

\begin{proof}
Write \(\ell_B=\log(B\operatorname{rd}_E)
 +|\RePart (L'/L)(1,\chi_K)|\), where \(\chi_K\) is the
quadratic character of the host construction. Assume the
determinant height bound \(\abs{H_{\tau,c}}\le C_f N\ell_{B_0}\)
on the prescribed component lines. The two calibrations in
Section~\ref{sec:covolume-flag-test} give
\[
 -a\le C(r_*)\ell_{B_0},\qquad |ja+Nb_0|\le C_f N\ell_{B_0},
\] with at most fifteen
level primes. Write \(m_\tau=N^{-1}\sum_i u_{\tau i}\).

For a decreasing list the centered first-block sum
\(b_\tau=\sum_{i<j}(u_{\tau i}-m_\tau)\) is nonnegative.
Also \(b_\tau\le\sum_i|u_{\tau i}|\) by writing it as the
first-block sum times \(1-j/N\) minus the second-block sum
times \(j/N\). The contribution of the shifts divided by
\(nVG(u)\) is \(\sum_\tau(a b_\tau+(ja+Nb_0)m_\tau)\).
The main term \(H(u)\) is nonnegative. Apply
Proposition~\ref{prop:flag-uniform-prime-tail} to the residual
prime terms in \eqref{eq:flag-fixed-shifts}. The lower bound on
\(a\) controls \(a b_\tau\) because \(b_\tau\ge0\), and the bound on
\(|ja+Nb_0|\) controls the other term since
\[
 N|m_\tau|\le\sum_i |u_{\tau i}|.
\]
Adding the estimates gives
\begin{equation}\label{eq:flag-slope-lower}
 \frac{\Phi_{B_0}(L_u^n)}{nVG(u)}
 \ge-C_*\left(\ell_{B_0}\sum_{\tau,i}|u_{\tau i}|
                           +j_s\sum_\tau U_\tau\right).
\end{equation}
One can use
\[
 C_*=C(r_*)+\max(C_f,0)+8\log4.
\]

In particular the constants from the wall and coefficient
extraction steps are absent.

\proofstep{An upper bound for the determinant height.}
On the marked model one can take \(C_f=11\).
At its interior test point in component \(c\), the empty-prime
determinant is the determinant of the \(j\) successive Lie
lines at \(\tau\) tensored with the inverse determinant of the
\(N-j\) lines at \(\bar\tau\). In a finite field of definition
take the two eigenframes on the semistable fibre. The local
determinant ideals are generated by their maximal minors in the
integral Lie lattice. Their Hermitian Gram determinants give
the squared complex norms. Exactness of the successive
filtrations makes the sum of the graded degrees the degree of
each determinant, including the order of each local ideal.
The flat-degree comparison for the polarized CM action identifies
their difference with
\[
 (2j-N)\RePart\log\Omega(E,\chi_K)/2.
\]

Here the comparison on full auxiliary level is applied to the
transported fibre and its two determinant models. Normalized
degree is invariant under its finite field extension. The CM
period estimate bounds the absolute value by
\(7N(\log\operatorname{rd}_E+|\RePart (L'/L)(1,\chi_K)|)\).
\label{cor:covolume-wall-shape-6}

Changing to the cutoff metric at the same point changes local
log norms only at the specified finite primes. The absolute
change per rational prime, after taking normalized place
weights, is at most
\[
 N\left(\frac{\log p}{p-1}+v_{\rm diff}(E,p)\log p\right).
\]
Indeed the local bound has the additional factor
\(-\log|p|_v/\log p\) and its normalized sum over \(v\mid p\)
is at most one in the conventions of the local reading.
The sum of the different terms is at most
\(\log\operatorname{rd}_E\). Summation up to \(B_0\)
with the exceptional primes included as in
Section~\ref{sec:covolume-flag-test} gives total change at most
\[
 4N(\log B_0+\log\operatorname{rd}_E).
\]
This proves the
claimed determinant bound on every component using the indicated
flat-degree comparison. \label{cor:covolume-wall-shape-7}

\proofstep{Passing to the component mean.}
For rational \(u\) the strict gaps imply that \(u-sd_0\)
is dominant for all sufficiently small positive rational
\(s\). Polynomial agreement gives limit \(P_{B_0,c}(u,0)\)
for the arithmetic self-intersection of \(M_{B_0}(u,s)\)
on \(c\). Its geometric degree tends to \(w_cG(u)>0\).
The quotient by \(n\) times the latter degree thus has
limit
\[
 s_c=P_{B_0,c}(u,0)/(nw_cG(u)).
\]
The geometric
component-volume comparison gives equal \(w_c\), hence
\[
   \frac{1}{\#\{c\}}\sum_c s_c=\frac{\Phi_{B_0}(L_u^n)}{nVG(u)} .
\]
This equality explains the occurrence of an arithmetic mean of
slopes on the deep cover. The bound in \eqref{eq:flag-slope-lower}
is for that mean. \label{cor:covolume-wall-shape-8}
\end{proof}
